\section{Introduction \texorpdfstring{\cred{[1.0 pages]}}{[1.0 pages]}}
\label{sec:intro}

As the world accelerates its adoption of agentic AI, it becomes crucial to build and assess the agents' capability and capacity in retaining information. 
In particular, to become a truly \textit{personal assistant} to offload tedious tasks, an AI agent needs to manage and act upon \textit{personal memory}, a collection of past perceptions, interactions, and experiences.
Edge-deployed agents such as OpenClaw~\citep{openclaw2026} and wearable memory systems such as Memoro~\citep{zulfikar2024memoro} make this setting real rather than hypothetical.

Consider this simple scenario: Bob asks Charlie’s personal agent assistant: “What did Alice say about her medical appointment?”

To answer correctly, the agent needs to perform a sequence of non-trivial tasks. It must search through weeks of daily conversations, find relevant mentions of the appointment, distinguish which parts of information shared by Alice are visible to Charlie,
%%XM \xm{What does this mean?} \jd{Former benchmarks are not ego-centric, so Charlie can see everyone's interaction, including all Alice's conversation. here i want to emphasize Charlie should not know about Alice's personal conversations}, 
and determine whether Alice later changed her mind. Then it must make an even harder judgment: even if the answer is known, is it actually appropriate to tell Bob?

This is the kind of challenge personal memory assistants will increasingly face.  The agent must continuously retain private interaction history, decide what to store, what to retrieve, what to compress, and what to withhold—all under the compute and privacy constraints of on-device deployment, where keeping sensitive interaction history local is part of the point. A benchmark for personal memory in this setting should therefore evaluate the \textit{full stack on device}, not just whether a model can retrieve the relevant information.

In addition to complexity, the sheer size of the data involved in such tasks is also easy to underestimate. 
Behavioral studies indicate that a person typically produces or hears on the order of ${\sim}10$K conversational tokens per day: starting from ${\sim}13$K--$16$K daily words, applying a social-conversation share 65\% and a 1.25--1.33 tokens/word conversion.\footnote{Using Mehl's ${\sim}$16K daily words and Tidwell's pooled mean of 12{,}792, then applying Dunbar's ${\sim}65\%$ social/personal share, yields ${\sim}$8.3K--10.4K conversational words/day~\citep{mehl2007women,tidwell2025talkative,dunbar1997gossip}. OpenAI, Gemini, and Anthropic English heuristics imply ${\sim}$1.25--1.33 tokens/word, \ie{} ${\sim}$10.4K--13.9K tokens/day, so we use ${\sim}$10K as a conservative round target~\citep{openai2026tokens,google2025geminitokens,anthropic2023claude2modelcard}.} After a few months, even one user's verbal history approaches the million-token scale before prompts, metadata, media data, or multi-party context are added. 
A million-token cloud context window does not help here, as it requires off-device upload, which privacy-sensitive personal memory management is meant to avoid. 
Instead, it falls on local, personal agents to effectively and safely maintain/utilize such dense human-human and human-agent interaction histories.

\para{Major structural gaps in personal memory evaluation.}
While existing long-context and memory benchmarks~\citep{maharana2024evaluating,wu2025longmemeval,memorybench2025} successfully evaluate recall in long transcripts, continuous learning, and memory-augmented assistants, they cannot assess the aforementioned \textit{personal memory capability/capacity}.
More specifically, there exist \underline{three structural gaps} between their design and our target use case.
\emph{(1)~Perspective}: existing benchmarks lump personal memory together, producing and building on an omniscient transcript of Alice, Bob, and others.
A true personal agent, such as Charlie's assistant, should only have visibility on conversations Charlie participated in or witnessed. 
(2)~\emph{Global coherence}: the context corpura used by most existing benchmarks are task-based, generated surrounding the actions and interactions involved in solving (often isolated) tasks. There is little or no temporal/spatial framework simulating human behavioral constraints (\eg, Alice and Bob are at the same location for a dinner discussion, and 70-year old Charlie would not send a message at 3am asking for bank transfer). Without global alignment in this regard, it is difficult to assess context coherence in evaluating information retrieval, reasoning, and protection. 
%%XM \st{the same people, events, locations, and access scopes recur across sessions in real deployment scenarios. Benchmarks that stitch independent fragments or noise can make retrieval and multi-session reasoning easier than deployment, because target evidence is often isolated in a topically separable fragment rather than entangled with recurring entities, later updates, and permission cues.}
\emph{(3)~Activity-density}: from the per-person/agent perspective, existing benchmarks provide very sparse data, generated through  a small number of isolated, task-based sessions. 
Evaluating at such a low density (far below the human-level typical daily volume given earlier) does not properly examine personal agents' persistent memory capacity, involving memory writes, interference, recency, and temporal provenance, all sensitive to dataset scale.

\para{Contributions.} We consider the major contributions of this work as follows: 
\begin{itemize}[leftmargin=1.5em,itemsep=-1pt,topsep=2pt]
  \item We propose \benchmark, the first ego-centric benchmark designed to evaluate \emph{personal on-device memory agents}, to our best knowledge.
  \item For corpus generation, \benchmark adopts \emph{world-grounded corpus generation} to produce synthetic human-human and human-agent interactions at typical human activity density. 
  To this end, we develop \masim, an evaluation-aware multi-agent simulator for generating long, coherent, temporally ordered personal-memory histories at scale. 
  \masim combines persona-driven social agents, person--person and person--agent dialog streams, provenance and access metadata, and day-batched synchronization: sessions within a day share a fixed memory snapshot and commit updates at day boundaries, preserving temporal consistency while enabling scalable corpus generation (\S\ref{sec:datagen}).
  \item For task generation, \benchmark adopts \emph{user-specific evidence projections}, which eliminates the need of manual annotation, and defines 6 task types spanning the recall, reasoning, and trustworthiness dimensions in accuracy (\S\ref{sec:benchmark}).
  Combining this with \masim, we create \benchL, a dense personal-memory dataset with 50 agents over 15 simulated days, 10.3M dialog-text tokens, around 24.1K text-only ego-observed tokens per agent per day, and evidence-linked evaluation instances (\S\ref{sec:setup}).
  \item Using \benchL, we evaluate five memory backends---Vanilla context, BM25-RAG~\citep{robertson2009bm25}, Oracle retrieval, Memobase~\citep{memobase2025}, and MemSearch---across five open-weight reader models, in both the accuracy dimensions named above and on-device latency/power consumption.
  Our results yield three major findings (\S\ref{sec:core-findings}, \S\ref{sec:ttft-finding}, \S\ref{sec:ablation-studies}): (i)~\emph{Two attractors on permission-aware access}: across the $25$ backend--reader cells, models cluster into ``privacy by amnesia'' (no evidence surfaced, low leak only because nothing is said) and ``anti-policy disclosure'' (evidence surfaced and disclosed despite the access marker, with the leak-minus-disclose gap growing monotonically from $+6$\,pp at Qwen3-0.6B to $+27$\,pp at Qwen3-32B); (ii)~\emph{Matched evidence dominates reader scale}: the smallest reader paired with Oracle context already outperforms every non-Oracle cell on Avg, and a $59$-cell retriever sweep plus an omniscient backend that hands the reader every session both stay $\geq 50$\,pp below the Oracle ceiling on cross-session reasoning, isolating matched-evidence provenance --- not retriever quality or context volume --- as the lever; (iii)~\emph{Search overhead bites only at the edge}: on Qwen3-0.6B the BM25 retrieval pipeline alone adds $87$\,ms to a $74$\,ms LLM-prefill, but on Qwen3-32B the same $87$\,ms is $<4\%$ of TTFT, so backend choice rather than reader choice drives perceived responsiveness only at the smallest reader.
\end{itemize}

% ============================================================================
% 2. RELATED WORK
